\documentclass[12pt]{article}

\usepackage[margin=1in]{geometry}
\usepackage{setspace}
\usepackage{natbib}
\setstretch{1.0}

\setlength{\parindent}{0pt}
\setlength{\parskip}{0.8em}

\title{This is a title}
\author{Wenqin Huang}
\date{\today}

\begin{document}

\maketitle


As low Earth orbit broadband networks expand rapidly, space launches have become increasingly indispensable to global connectivity, yet their growing environmental costs can no longer be treated as a secondary concern. Space activity is no longer occasional; it is becoming a rapidly expanding industrial sector: according to Ryan et al.'s 2022 research, the scale of it will leap from 350 billion in 2019 to over a trillion by 2040~\citep{ryan2022ozone,bushnell2018newspace}. This critical leap will probably cause a substantial intensification and amplification of its environmental impacts. Additionally, by rocket launching activities, multiple pollutants will be directly injected into the upper atmosphere, where pollution will be far more destructive, exposing a direct threat to ozone~\citep{ryan2022ozone,kokkinakis2022pollution}. Moreover, this expansion is not only what will happen in the future, but already ongoing: the frequency of rocket launching increased substantially from 58 times in 2003 to 329 times in 2025~\citep{mcdowell2025space}. The pollution is not simply limited in the earth ecosystem, but also influences the earth orbits. Frequent satellite constellations will also occupy earth orbits, create hundreds of thousands of pieces of space debris, and subsequently obstruct future satellite deployments~\citep{torres2020reusable}. Thus, the environmental impact of rocket launching could no longer be ignored, since the industry is rapidly expanding, and it influences the environment in multiple different ways.



However, the constellation of LEO broadband is not simply a symbolic space activity, but an important component of communication infrastructure. The rapid expansion of OneWeb and Starlink is motivated and triggered by the growth of global communication demand. The connectivity created by LEO broadband systems is also widely argued to be beneficial, especially in remote areas, where traditional internet doesn't cover: LEO satellites provide them internet connection, and therefore improve their social connection, and subsequently mental health \citep{osoro2025leo,sandau2010smallsat}. Based on the benefit it brings, space launches should not simply be halted. Thus, the true essential problem is "to what extent can technological strategies reduce these environmental costs and improve the sustainability", rather than "whether people should continue launching or not".



To address the environmental costs associated with frequent space launches, various technological strategies have been proposed to improve the sustainability of launch systems. Among these, reusable launch vehicles, alternative propellants, and satellite miniaturization are often considered to be the most practical approaches. Their effects can be evaluated based on their ability to reduce emissions, minimize life-cycle environmental impacts, and remain effective under the rapid expansion of satellite deployment. However, the effectiveness and validity of these approaches are not essentially limited by technological capability, but are complicated by financial incentives, technological limitations, and the rebound effect that may offset or even reverse their environmental benefits. Therefore, although technological strategies are already sufficiently developed to reduce the environmental impact of space launches, their real-world effectiveness remains constrained by economic barriers, industrial inertia, and the rapid expansion of LEO megaconstellations, limiting their practical implementation.



Alternative propellants are the most direct technique to reduce environmental impact among three approaches, since they change the fuel combustion product directly. Dallas et al. systematically reviewed the environmental impact brought by traditional propellants, and argue that solid fuel, which typically involves chemicals like ammonium perchlorate, may potentially damage the ozone the most; kerosene-based propellants mainly emit carbon dioxide (CO$_2$) and black carbon, which may intensify the greenhouse effect, but normally are less destructive than solid propellants~\citep{dallas2020review,brown2024inventory}. By contrast, alternative fuels are suggested to be more environmentally friendly. Liquid hydrogen is widely recommended as a candidate for future propulsion, since the emission it generates contains mainly water vapor, and with negligible ecological toxicity~\citep{dallas2020review,barato2023fuels}. To be more specific, Calabuig et al.'s study suggests that the carbon footprint of liquid hydrogen (alternative propellant) fleets is typically 2--8 times less than methane-based (transitional propellant, cleaner than kerosene, but not as clean as hydrogen) vehicle fleets. Additionally, they further argue that, in most cases, choosing sustainable propellants will affect more than the vehicle's reusability: propellant choice is the main differentiation factor when it's under environmental criteria~\citep{calabuig2024reusable}. Recent research by Hai et al. suggests that environmentally friendly propellants not only reduce toxic products, but also lower the risk of transporting and operating~\citep{hai2026propellants}. These findings collectively suggest that alternative propellants provide a highly effective approach for reducing the environmental impact of rocket launches in principle. However, their real-world implementation remains obstructed by practical and economic challenges. The storage and transportation for alternative propellants are typically complex and expensive. Dallas et al. argue that liquid hydrogen is inherently unsuitable for storage: its density is lower, which means it requires a storing system that is significantly larger and more complex than traditional fuels; additionally, the boiling point of liquid hydrogen is considerably lower, which means the propellant will be more likely to vaporize---existing propellant storing tanks may not be suitable for storage~\citep{dallas2020review}. This indicates that the cleaner propellant may not be easier to store or use. Additionally, Calabuig et al. also argue that although the technology is proven beneficial, the cost is still significantly higher than traditional, and the industry, including producing, transporting, and storing, is also immature. Since more and more rocket launching activities are nowadays carried out by private sectors, rather than government-supported institutions, the higher cost may discourage companies from utilizing alternative propellants~\citep{calabuig2024reusable}. Therefore, the established aerospace industry was structured around traditional propellants like kerosene, and the transformation may not be financially favorable. In conclusion, alternative propellants offer a concrete environmentally friendly choice and emission reduction strategy, but their effect depends on the pace of industrial transformation and financial incentives, rather than technological limitations.



However, propellant choice alone does not fully determine the environmental impact of space launches. Beyond fuel composition, the design and operational characteristics of launch systems---particularly the use of reusable launch vehicles---also play a critical role. Reusable launch vehicles offer a promising solution to reduce the environmental costs of space launches by reducing the manufacturing emissions and wastes. Therefore, applying reusable rocket technology can reduce the manufacturing cost for launching, but it will not necessarily reduce net environmental impact, Calabuig et al. argue~\citep{calabuig2024reusable}. On one hand, reusable vehicles can improve the material efficiency of fixed equipment, the fleet, and therefore reduce the launching cost financially. As Calabuig writes, the original idea of reusable rockets is to reduce material wasting, and when evaluating under the criterion of launching impact per unit of mass, reusable vehicles can indeed improve both environmental and financial efficiency. However, Calabuig et al. also provide solid numerical evidence proving that the improvement by reusable rockets is limited: for liquid hydrogen fleets, the climate impact could not be reduced over 40\%, for methane fleets, the impact is less significant, around 20\%~\citep{calabuig2024reusable}. This trend indicates that the climate impact reduction is considerable, but not as transformative as alternative propellants. However, long-term assessments, like life-cycle evaluations, are also critical in evaluating environmental impact of an industry~\citep{wilson2022lifecycle}. While under long-term evaluation, because such vehicles are more financially favorable, with same budgets, companies may launch substantially more frequently and combust a significantly greater amount of propellants, which, as mentioned in last paragraph, could amplify the environmental impact and carbon emissions to a greater extent, causing a rebound effect. The rebound effect of reusable vehicles is well illustrated by SpaceX, a private launch provider that has pioneered the large-scale adoption of reusable rocket technology, particularly through its Falcon 9 program. According to FAA environmental reviews, the annual Falcon launch cadence considered for SpaceX has risen steadily, from 20 launches per year at LC-39A and 50 at SLC-40 in the 2020 assessment, to 36 launches per year at Vandenberg in 2023, 50 in 2024, and 100 in the 2025 Vandenberg EIS \citep{faa2025falconprogram}, while it reaches 165 times of launching in 2025, this single year~\citep{mcdowell2025space}. These statistics confirm the launching frequency's leap under the implementation of reusable vehicle fleets. To summarize, although reusable launch vehicles can improve efficiency by reducing manufacturing impacts and lowering costs, their overall environmental benefits remain limited and could be offset, or even reversed, by increased launch frequency driven by economic incentives.



Satellite miniaturization improves the functional efficiency of individual satellites, but it also lowers the barrier to large-scale deployment. As Sandau et al. suggested, satellite miniaturization could improve the efficiency per unit of mass by reducing the mass while the function remains unchanged, and enable the satellite broadband system to cover globally~\citep{sandau2010smallsat}. Satellite miniaturization also makes satellites cheaper, easier to design, and to deploy. For environmental impact, the miniaturization of satellites means that, with same rocket, more satellites could be launched and deployed, therefore, the average environmental cost of each satellite could be reduced. However, the miniaturization of satellites could also increase the scale of broadband constellation, while the increase of scale may not only intensify environmental burden, but also occupy the orbit and obstruct future launches~\citep{kumaran2024manufacturing}. While megaconstellations enabled by miniaturization significantly improve broadband access in remote and underserved regions, their emissions intensity per subscriber remains substantially higher than that of terrestrial 4G systems, indicating a non-negligible environmental cost~\citep{osoro2025leo,kukreja2025leo}. This indicates that the benefits of satellite miniaturization are highly context-dependent: although efficiency is improved at the level of individual satellites, the expansion of constellation scale shifts the problem from per-unit optimization to total system impact.



After analyzing three technological strategies including alternative propellants, reusable rockets and satellite miniaturization, we can find that all of three technologies are theoretically capable of reducing the environmental impact of satellite megaconstellations, but their effectiveness in real context is still largely limited. The key to addressing this problem is not about the effectiveness of technology, but depends on the complicated economic and industrial environment. Firstly, the market's incentive mechanism wasn't designed for optimizing environmental goals: the reduction of cost could not reduce total launches, but trigger more frequent spaceflight missions and deployments. To be more specific, the development of reusable rockets escalates spaceflight frequency by reducing its cost, and satellite miniaturization expands the megaconstellation of broadband networks by lowering per-satellite manufacturing cost. This scope expansion driven by technological advancements and efficiency improvements makes their net environmental benefits offset, and even trigger reverse effects. Secondly, the existing aerospace industry exhibits significant path dependence. For long, the rocket industry, including launchpad infrastructures, propellant storage systems, and related supply chains, has been established around traditional fuel systems. This means that although alternative propellants may offer better environmental performance, they may still face limitations in implementation and transformation due to the high infrastructure-reconstruction cost and industrial inertia. This inertia makes it hard to practically transfer this technological effectiveness to large-scale adoption. Furthermore, to a great extent, the environmental impacts are derived from the expansion of scale, rather than solely caused by technological deficiency and inefficiency. In this case, it will be difficult to control the environmental burden solely relying on technological advancements and optimizations. Therefore, these evidence and resonings collectively suggest that the main factor limiting spaceflights' sustainability is implementation and transformation, rather than technological limitations. This also means that, in order to make aerospace activities sustainable, institutional frameworks and policy tools, like limiting launching frequency and guiding industrial transition, are also required.



In conclusion, the rapid expansion of LEO broadband satellite constellations has intensified the environmental challenges associated with space launches, making sustainability a critical concern for the future of the space industry. While technological advancements, like alternative propellant, reusable rockets and satellite miniaturization, offer practical approaches to reduce environmental impacts, their effectiveness remains controvertial and limited in practice. Alternative propellants provide the most direct way for reducing emission, but face financial barriers and industrial inertia. Reusable launch vehicles improve per-launch efficiency and reduce cost financially, but their benefits are limited by the rebound effect driven by cost reductions. Similarly, satellite miniaturization enhances functional efficiency, but often leads to large-scale deployment that increases total environmental burden.



These findings suggest that the primary limitation is not the lack of technological capability, but the implementation of these existing Technologies. As long as economic incentives favor expansion and existing industrial systems doesn't transit, technological improvements alone are not quite possible to achieve significant reductions in environmental impact. Therefore, ensuring the sustainability of space activities requires not only technological innovation, but also structural changes in political regulatory frameworks and economic incentives. Finally, sustainable space development depends on promoting technological progresses' application with system-level regulation, rather than relying on methodological advancements alone.

%\subsection{Why Mature Technologies Remain Underutilized}

%\subsection{Barriers to Adoption}

%\subsection{The Case for Regulatory Intervention}

%\section{Conclusion}

\newpage
\bibliographystyle{apalike}
\bibliography{references}


\end{document}
